\documentclass{article}
\usepackage{graphicx} % Required for inserting images
\usepackage{xspace}
\usepackage[utf8]{inputenc}
\usepackage[T1]{fontenc}
\usepackage[naustrian]{babel}
\usepackage{amsmath}

\title{Zur Vorteilhaftigkeit direkter adiabater Kühlung in Wohnräumen}
\author{}
\date{}

\begin{document}

\maketitle

\section{Einleitung}

Der Hitzeindex $HI$ ist eine Zustandsgröße und kann wie folgt beschrieben werden:

\begin{equation}
    HI = c_1 + c_2 T + c_3 \varphi + c_4 T \varphi + c_5 T^2 + c_6 \varphi^2 + c_7 T^2\varphi + c_8 T \varphi^2 + c_9 T^2 \varphi^2
    \label{eq:hitzeindex}
\end{equation}

mit der Umgebungstemperatur $T$ in Grad Celsius und der relativen Luftfeuchtigkeit $\varphi$ in Prozent.
Die Intention des Hitzeindex ist ein Maß, welches das menschliche Wärmeempfinden besser darstellt als die Trockenkugeltemperatur.

\subsection{Luftfeuchtigkeit}
Die relative Luftfeuchtigkeit beträgt
\begin{equation}
    \varphi = \frac{\rho_w}{\rho_{w, max}}
    \label{eq:relfeuchte}
\end{equation}

Der Zusammenhang zwischen der Sättigungsmenge $\rho_{w, max}$ in g\,m$^{-3}$ und der Temperatur in Grad Celsius $T$ lässt sich durch folgende Gleichung beschreiben:
\begin{equation}
    \rho_{w, max} = 6,0495 \cdot e^{0,0533 \cdot T}
    \label{eq:sättigung}
\end{equation}

Zurück eingesetzt in Gleichung \ref{eq:relfeuchte}:
\begin{equation}
    \varphi = \frac{\rho_w}{6,0495 \cdot e^{0,0533 \cdot T}} = 0,16530 \cdot \rho_w \cdot e^{-0,0533 \cdot T}
\end{equation}

Und zurück eingesetzt in Gleichung \ref{eq:hitzeindex}:
\begin{multline}
    HI = -8,784695 + 1,61139411 \cdot T + 0,3865 \cdot \rho_w \cdot e^{-0,0533 \cdot T} - \\ - 0,024153 \cdot T \cdot \rho_w \cdot e^{-0,0533 \cdot T} - 0,01231 \cdot T^2 - 0,002715 \cdot \rho_w \cdot e^{-0,1066 \cdot T} + \\ + 0,0003656 \cdot T^2 \cdot \rho_w \cdot e^{-0,0533 \cdot T} + 0,0001199 \cdot T \cdot \rho_w \cdot e^{-0,1066 \cdot T} - \\ - 5,921 \cdot 10^{-7} \cdot T^2 \cdot \rho_w \cdot e^{-0,1066 \cdot T}
\end{multline}

Leiten wir den Hitzeindex einfach mal zur Gaudi nach der Temperatur ab:
\begin{multline}
    \frac{\partial HI}{\partial T} = 6.311786 \cdot 10^{-8}\, \rho_w\, e^{-0.1066\, T}\, T^2 - \\ - 1.94865 \cdot 10^{-5}\, \rho_w\, e^{-0.0533\, T}\, T^2 - 1.39655 \cdot 10^{-5}\, \rho_w\, e^{-0.1066\, T}\, T + \\  + 0.00201855\, \rho_w\, e^{-0.0533\, T}\, T + 0.000409319\, \rho_w\, e^{-0.1066\, T} - \\ - 0.0447535\, \rho_w\, e^{-0.0533\, T} - 0.02462\, T + 1.61139
\end{multline}

Mir war das diese Schreibübung händisch zu blöd, darum an dieser Stelle Dank an meinen Coautor!

Und weil es so lustig war, das Ganze nochmal mit der Luftfeuchte:
\begin{multline}
    \frac{\partial HI}{\partial \rho_w} = -5.921 \cdot 10^{-7}\, e^{-0.1066\, T}\, T^2 + 0.0003656\, e^{-0.0533\, T}\, T^2 + \\ + 0.0001199\, e^{-0.1066\, T}\, T - 0.024153\, e^{-0.0533\, T}\, T - \\ - 0.002715\, e^{-0.1066\, T} + 0.3865\, e^{-0.0533\, T}
\end{multline}

Wer hätte gedacht, was jetzt kommt?
\begin{equation}
    \mathrm{d} HI = \left( \frac{\partial HI}{\partial T} \right)_{\rho_w} \mathrm d T + \left( \frac{\partial HI}{\partial \rho_w} \right)_T \mathrm d \rho_w
\end{equation}

\subsection{Thermodynamik}
Die Verdampfungsenthalpie von Wasser beträgt
\begin{equation}
    \Delta h_v = 2441 \, \mathrm{kJ} \cdot \mathrm{kg}^{-1}
\end{equation}
während die Wärmekapazität von Wasser eh niemanden interessiert.

Die Wärmekapazität von Luft beträgt etwa
\begin{equation}
    c_p = 1,01 \, \mathrm{kJ} \cdot \mathrm{kg}^{-1} \cdot \mathrm{K}^{-1}
\end{equation}

Der erste Hauptsatz der Thermodynamik lautet für unser System:
\begin{eqnarray}
    \mathrm d h &=& \mathrm d q - p \cdot \mathrm d V \\
    \mathrm d h &=& 0 \\
    0 &=& c_p \mathrm d T + \Delta H_v \mathrm d x \\
    \Delta H_v \mathrm d x &=& - c_p \mathrm d T \\
    \frac{\mathrm d x}{\mathrm d T} &=& -\frac{c_p}{\Delta H_v} \\
    \frac{\mathrm d x}{\mathrm d T} &=& -0,41 \, \mathrm K \cdot \mathrm{g}_w \cdot \mathrm{kg}_L^{-1}
\end{eqnarray}

mit der Luftdichte $\varrho_L = 1,2041 \,$ kg m$^{-3}$ (definitiv konstant und nicht abhängig von Temperatur oder Luftfeuchtigkeit) ergibt sich:
\begin{equation}
    \frac{\mathrm d \rho_w}{\mathrm d T} = -0,34 \, \mathrm K \cdot \mathrm{g}_w \cdot \mathrm{m}_L^{-3}
\end{equation}

\section{Differentialgleichung}
\begin{eqnarray}
    \mathrm{d} HI &=& \left( \frac{\partial HI}{\partial T} \right)_{\rho_w} \mathrm d T + \left( \frac{\partial HI}{\partial \rho_w} \right)_T \mathrm d \rho_w \\
    \frac{\mathrm d HI}{\mathrm{d T}} &=& \left( \frac{\partial HI}{\partial T} \right) + \left( \frac{\partial HI}{\partial \rho_w} \right) \frac{\mathrm d \rho_w}{\mathrm d T}
\end{eqnarray}

Wenn wir einsetzen, kommt dieses unlesbare Monstrum raus:

\begin{multline}
    \mathrm d HI = ( 6.311786 \cdot 10^{-8}\, \rho_w\, e^{-0.1066\, T}\, T^2 - \\ - 1.94865 \cdot 10^{-5}\, \rho_w\, e^{-0.0533\, T}\, T^2 - 1.39655 \cdot 10^{-5}\, \rho_w\, e^{-0.1066\, T}\, T + \\  + 0.00201855\, \rho_w\, e^{-0.0533\, T}\, T + 0.000409319\, \rho_w\, e^{-0.1066\, T} - \\ - 0.0447535\, \rho_w\, e^{-0.0533\, T} - 0.02462\, T + 1.61139 ) + \\ + (2.01 \cdot 10^{-7}\, e^{0,036\, T}\, T^2 - 0,00012\, e^{-0.0533\, T}\, T^2 - \\ - 0,000041\, e^{-0.1066\, T}\, T + 0,0082\, e^{-0.0533\, T}\, T + \\ + 0,00092\, e^{-0.1066\, T} - 0,13\, e^{-0.0533\, T})
\end{multline}

\end{document}
